% Image credits for the photographs and AI illustrations of Book 2
% (University Chemistry, Year 1), Dutch edition. Machine-readable license
% records live in images/book2/CREDITS.md; keep this file in sync with
% frontmatter/image-credits-book2.tex.
\chapter*{Beeldverantwoording}

De tekeningen in dit boek zijn origineel. De foto's worden gebruikt onder
hun respectieve vrije licenties, met dank aan hun makers; de volledige
bronlinks en licentie-URL's van elke foto staan in
\texttt{images/book2/CREDITS.md} in de repository van het boek.
De gevarenpictogrammen zijn die van het Globally Harmonized System of
Classification and Labelling of Chemicals, zoals getekend voor het pakket
\texttt{ghsystem} van \TeX\ Live.

\bigskip
Naast de foto's en de originele tekeningen zijn sommige figuren
illustraties die met AI gegenereerd zijn, gemaakt met de beeldgeneratie
van OpenAI op basis van prompts die door de redactie van het boek
geschreven zijn, en vóór opname nagekeken op chemische juistheid. De
volledige prompt van elk gegenereerd beeld staat in
\texttt{images/book2/ai/PROMPTS.md} in de repository.

\bigskip
\noindent Foto's, per hoofdstuk:
\begin{itemize}\small
  \item Hfst.~1: Niels Bohr, rond 1922, AB Lagrelius \& Westphal, publiek
        domein (Wikimedia Commons).
  \item Hfst.~2: Linus Pauling, 1962, Nobel Foundation, publiek domein
        (Wikimedia Commons).
  \item Hfst.~5: gedegen koper, door Flipin, OG, CC0 (Wikimedia Commons).
  \item Hfst.~6: halietkristallen, fluoriet en een octaëdrische diamant,
        door James St. John, CC BY 2.0; grafiet, door Zbynek Burival,
        CC BY-SA 4.0; sneeuwkristallen door Wilson A. Bentley (1902),
        publiek domein (alle Wikimedia Commons).
  \item Hfst.~8: Svante Arrhenius, 1909, Meisenbach Riffarth \& Co.,
        publiek domein (Wikimedia Commons).
  \item Hfst.~11: neerslagen van zilverchloride, -bromide en -jodide, door
        Rrausch1974, CC BY-SA 3.0 (Wikimedia Commons).
  \item Hfst.~12: suspensie van koper(II)hydroxide, door Alvy16, CC BY 4.0;
        oplossing van een koperamminecomplex, door Chemicalinterest,
        publiek domein (beide Wikimedia Commons).
  \item Hfst.~13: Walther Nernst, 1889, Smithsonian Libraries, publiek
        domein (Wikimedia Commons).
  \item Hfst.~16: Jacobus Henricus van 't Hoff, foto gepubliceerd in 1911,
        publiek domein (Wikimedia Commons).
  \item Hfst.~21: Victor Grignard, publiek domein (Wikimedia Commons).
  \item Hfst.~25: Vladimir Markovnikov, portret uit een in memoriam
        gepubliceerd in 1905, publiek domein (Wikimedia Commons).
  \item Hfst.~26: lithiumbekkens van de Salar de Atacama, NASA Earth
        Observatory, publiek domein; natrium in minerale olie, door Justin
        Urgitis, CC BY-SA 2.5; vlamkleuren, door Hegelrast, CC BY-SA 4.0
        (alle Wikimedia Commons).
  \item Hfst.~27: pisolitische bauxiet, door James St. John, CC BY 2.0;
        borax, door Aram Dulyan, publiek domein; kwarts, door Stephanie
        Clifford, CC BY 2.0; witte fosfor, Hi-Res Images of Chemical
        Elements, CC BY 3.0; Fritz Haber, The Nobel Foundation, publiek
        domein (alle Wikimedia Commons).
  \item Hfst.~28: zwavelwinning in de krater van de Kawah Ijen, door
        S\'emhur, CC BY-SA 4.0; chloor in een ampul, door W. Oelen,
        CC BY-SA 3.0; broom in een ampul, door Jurii, CC BY 3.0;
        joodkristallen, door Dnn87, CC BY 3.0 (alle Wikimedia Commons).
  \item Hfst.~29: verwarmde bank (koflerbank), door HBR, CC BY-SA 3.0;
        abberefractometer, door Foreade, CC BY-SA 4.0 (beide Wikimedia
        Commons).
\end{itemize}
\clearpage
